\section*{Limitations}

\textsc{RetroGen} relies on the availability of high-quality expert artifacts. While such artifacts are often more abundant than process annotations, their usefulness depends on whether they contain enough recoverable evidence and structural signals to support retrospective reconstruction. In domains where final outputs are highly underspecified, stylistically diverse, or weakly grounded in explicit evidence, the reconstructed trajectories may be less reliable.

Our experiments focus on evidence-grounded long-form generation with retrieval-oriented tools. Although the framework is designed to be general, further work is needed to evaluate retrospective process supervision in domains requiring richer interactive environments, longer-horizon planning, or non-textual tools such as code execution, database operations, and multimodal perception.